\chapter{End-to-End Reconstruction of a 4K Capture}
\label{ch:reconstruction}

This chapter presents a complete, worked case study in which a single hand-held 4K gallery walkthrough video was processed through the \texttt{gaussian-toolkit} pipeline — from raw footage to a trained, evaluated Gaussian field — without manual intervention beyond initiating the job. The account proceeds stage by stage, covering every transformation the data undergoes, with measured outputs reported at each step. Downstream tasks, specifically object decomposition and exportable mesh generation, are treated in Chapters~\ref{ch:scenegraph} and~\ref{ch:meshing} respectively. The objective of the present chapter is to establish the baseline volumetric reconstruction quality and the computational cost upon which all subsequent processing stages build.

The case study demonstrates a key design principle of the \texttt{gaussian-toolkit}: that the entire reconstruction pathway — frame selection, \acrshort{sfm}, and Gaussian training — is driven by a single pipeline invocation operating on a single video file, with no intermediate user decisions, no manual tie-point marking, and no specialised capture equipment. The source footage was collected with a consumer smartphone during a routine gallery visit. The resulting Gaussian field is photorealistic, geometrically coherent, and renderable in real time.

% ──────────────────────────────────────────────────────────────
\section{The Source Capture}
\label{sec:recon-capture}

The source material is a single, continuous video file recorded on an iPhone~15~Pro at 3840\,×\,2160 pixels (4K~UHD), 30~frames per second, with optical image stabilisation active and the main 24\,mm-equivalent wide camera selected throughout. The recording spans approximately five minutes of walking time, yielding a raw file of approximately 9\,GB in the Apple ProRes~4444 container format. The scene is a gallery interior housing multiple discrete artworks: five oil paintings in carved gilt frames hung on painted masonry walls, two freestanding bronze sculptures on white plinth pedestals, and a backlit cabinet installation containing archaeological replica artefacts.

The operator walked a loose figure-of-eight path through the space, pausing for approximately eight seconds in front of each principal artwork to allow multi-angle coverage, and sweeping the camera through wide arcs during transitions between works to ensure the pipeline could recover continuous camera trajectories. No tripod, motorised slider, or drone was used. The capture protocol required no specialist skill beyond keeping the scene in frame; it is replicable by any gallery staff member with a contemporary smartphone.

\begin{figure}[htbp]
  \centering
  \begin{subfigure}[t]{0.48\textwidth}
    \centering
    \includegraphics[width=\textwidth]{v2/src_frame.jpg}
    \caption{Frame at approximately 0\,min\,42\,s — central gallery bay, showing the spatial relationship between the sculpture pedestals and the hung paintings beyond.}
    \label{fig:src-frame-a}
  \end{subfigure}
  \hfill
  \begin{subfigure}[t]{0.48\textwidth}
    \centering
    \includegraphics[width=\textwidth]{v2/src_frame_b.jpg}
    \caption{Frame at approximately 3\,min\,17\,s — close approach to a framed oil painting, with specular highlights on varnish visible at the painting centre.}
    \label{fig:src-frame-b}
  \end{subfigure}
  \caption{Representative source frames from the iPhone~15~Pro 4K capture. The mixed tungsten--fluorescent--natural-light environment, specular highlights on varnished surfaces, and shallow depth-of-field transitions in the close-up frame are all characteristics that the pipeline must tolerate.}
  \label{fig:src-frames}
\end{figure}

Three characteristics of hand-held smartphone capture are worth examining in relation to pipeline tolerance, as they represent the most common failure modes in consumer-video reconstruction workflows.

\paragraph{Rolling-shutter distortion.} The iPhone~15~Pro CMOS sensor reads out line-by-line rather than all at once, causing moving subjects to appear sheared during rapid panning motions. The \texttt{gaussian-toolkit} does not apply an analytical rolling-shutter correction, which would require known pixel-readout timing. Instead, the Fibonacci-sphere keyframe selector (Section~\ref{sec:recon-keyframe}) penalises frames with high inter-frame optical-flow magnitude, which correlates strongly with rolling-shutter artefact severity. Frames extracted during slow, deliberate camera movements — the dominant motion mode in gallery walkthroughs — are unaffected, and the high-flow transition frames are filtered before reaching \acrshort{colmap}.

\paragraph{Motion blur.} Low-illuminance transition corridors between gallery bays caused brief periods of motion blur at the corners of the figure-of-eight path. The Laplacian-variance sharpness score embedded in the keyframe scoring function (Equation~\ref{eq:fib-score}) is specifically designed to detect such blur: the variance of the Laplacian of a blurred image is substantially lower than that of a sharp image at the same resolution. Blurred frames receive near-zero sharpness scores and are therefore excluded from the selected keyframe set irrespective of their angular coverage contribution.

\paragraph{Exposure variation.} The gallery's lighting mix — tungsten spotlights focused on the paintings, ambient fluorescent ceiling strips, and natural daylight admitted through a south-facing roof lantern — created per-frame auto-exposure swings of roughly $\pm$1.5~stops across the five-minute sequence. Rather than normalising exposures during pre-processing, the pipeline relies on the per-camera exposure compensation terms that are jointly optimised during \acrshort{3dgs} training \parencite{Kerbl2023}. Each registered camera is assigned a learnable scalar that scales the photometric loss for that viewpoint, allowing the optimiser to account for lighting variation without corrupting the scene's true appearance model.

% ──────────────────────────────────────────────────────────────
\section{Frame Extraction and Keyframe Selection}
\label{sec:recon-keyframe}

The raw \texttt{.mov} file was ingested by the pipeline's first processing stage, which invoked FFmpeg to decode and extract frames at a fixed temporal stride of one frame every 0.5~seconds — equivalent to 2~frames per second from the 30\,fps source — yielding \textbf{600~candidate frames} at full 4K resolution. The choice of 2\,fps represents a deliberate engineering trade-off: at the walking pace observed in the capture (approximately 0.8\,m/s at the camera), a 0.5-second stride corresponds to roughly 40\,cm of lateral displacement, which produces sufficient inter-frame parallax for \acrshort{sfm} feature matching whilst avoiding the computational burden of processing all 9,000 source frames. Experiments with 1\,fps decimation on similar indoor captures produced noticeably gapped angular coverage; 3\,fps provided negligible quality improvement over 2\,fps whilst increasing \acrshort{sfm} runtime by approximately 80\,\%.

From the 600 candidates, a subset of 240~keyframes was selected by the \texttt{fibonacci\_sampler.py} module, which implements a viewpoint-diversity scoring scheme based on a discretised Fibonacci-sphere angular grid. The Fibonacci sphere is an efficient quasi-uniform sampling of the unit sphere that avoids the clustering artefacts of latitude--longitude grids; projecting estimated camera orientations onto it allows the selector to measure how well the current frame set covers the full hemisphere of possible viewing directions. Each candidate frame $i$ is assigned a composite score:

\begin{equation}
  s_i = 0.6\,\sigma_i + 0.4\,\kappa_i
  \label{eq:fib-score}
\end{equation}

\noindent where $\sigma_i \in [0,1]$ is the Laplacian-variance sharpness score, normalised to the empirical maximum across all candidates, and $\kappa_i \in [0,1]$ is the angular coverage score, computed as the inverse of the density of previously selected viewpoints in the Fibonacci cell nearest to frame~$i$'s estimated camera direction. Frames are ranked by $s_i$ and selected greedily in descending order, subject to a minimum angular separation constraint of $5°$ between any two selected frames to prevent tight clusters of near-identical viewpoints. The relative weights of 0.6 (sharpness) and 0.4 (coverage) were determined empirically on a held-out set of indoor walkthrough sequences and reflect the observation that a sharp frame in a well-covered region is preferable to a blurred frame in an under-covered region, but that coverage is not negligible.

\begin{table}[htbp]
  \centering
  \caption{Frame selection pipeline summary for the gallery capture.}
  \label{tab:frame-selection}
  \begin{tabular}{@{}L{7.0cm} R{3.0cm}@{}}
    \toprule
    \textbf{Stage} & \textbf{Frame count} \\
    \midrule
    Source video (5\,min at 30\,fps)              & 9,000 \\
    Candidates after 2\,fps temporal decimation   & 600 \\
    Keyframes selected (Fibonacci scoring)        & 240 \\
    \bottomrule
  \end{tabular}
\end{table}

The reduction from 600 candidates to 240 keyframes — a 60\,\% reduction — translates directly into faster \acrshort{sfm} feature matching and shorter Gaussian training runs, because both stages scale with input image count. Critically, the reduction does not sacrifice viewpoint diversity: the coverage term $\kappa_i$ in the scoring function ensures that undersampled angular regions continue to attract new keyframes until the selection budget is exhausted. A post-selection analysis of the 240 keyframes confirmed coverage of 94\,\% of Fibonacci cells in the upper hemisphere, with the only unsampled cells corresponding to downward-looking angles that the operator did not achieve during the floor-level walkthrough.

% ──────────────────────────────────────────────────────────────
\section{Structure-from-Motion}
\label{sec:recon-sfm}

The 240 selected keyframes were passed to \acrshort{colmap} \parencite{Schoenberger2016sfm,Schoenberger2016mvs} configured in sequential exhaustive matching mode. COLMAP first extracted SIFT keypoints and descriptors from each image at the full 4K resolution, then matched descriptors across all image pairs within a sliding temporal window of $\pm$30 frames to balance recall against runtime. Camera intrinsics were modelled using the \texttt{OPENCV} distortion model (radial and tangential terms), which correctly represents the iPhone~15~Pro's primary lens distortion profile. Initial experiments with the simpler \texttt{SIMPLE\_RADIAL} model produced mean reprojection residuals approximately 0.2\,px higher and introduced systematic curl distortions in the recovered point cloud, consistent with under-modelled tangential distortion. The \texttt{OPENCV} model was therefore used throughout.

Following feature extraction and matching, COLMAP's incremental \acrshort{sfm} registered cameras sequentially, seeding the reconstruction from the image pair with the highest number of verified correspondences and extending outward. Bundle adjustment was applied globally at regular intervals and once more at the conclusion of reconstruction to minimise reprojection error across all registered cameras and 3D points jointly.

The final reconstruction registered \textbf{232 of the 240 input images} (96.7\,\%), producing a sparse point cloud of approximately \textbf{210,000 3D points} with a mean reprojection error of \textbf{0.64\,px}. This figure is well within the accepted threshold of 1.0\,px for high-quality \acrshort{sfm} initialisation of Gaussian splat training, and compares favourably with typical reprojection errors reported for indoor photogrammetric surveys conducted with dedicated mirrorless cameras.

The 8 unregistered frames (3.3\,\%) arise from two distinct failure modes. Five frames were captured during rapid camera rotations between gallery bays, where the angular displacement between consecutive frames exceeded the SIFT matching window despite passing the sharpness filter — a known boundary condition for sequential matching when the frame sequence includes abrupt direction changes. The remaining three frames show the interior of the backlit cabinet installation, where the combination of strong transmitted light and highly reflective glass surfaces destabilised SIFT descriptors by producing inconsistent keypoint responses across viewing angles. These 8 frames were removed from the input set automatically; because the keyframe selector had ensured broad spatial coverage, no visible reconstruction gaps resulted from their omission.

\begin{technote}
  COLMAP's intrinsic model selection has a measurable effect on downstream Gaussian training quality. In the gallery capture, switching from \texttt{SIMPLE\_RADIAL} to \texttt{OPENCV} reduced the mean reprojection error by 0.2\,px and increased the number of registered cameras from 228 to 232. The four additional cameras were all acquired at oblique angles to the cabinet installation where tangential distortion is most pronounced. The improvement propagated through to a 0.3\,dB gain in the trained model's \acrshort{psnr}, confirming that intrinsic model fidelity matters for reconstruction quality even with modern Gaussian splat optimisers that apply per-camera compensation terms.
\end{technote}

% ──────────────────────────────────────────────────────────────
\section{Gaussian Splat Training}
\label{sec:recon-training}

The COLMAP sparse point cloud and registered camera poses were imported directly into LichtFeld Studio \parencite{LichtFeldStudio2026}, an open, real-time \acrshort{3dgs} training and editing suite. LichtFeld Studio extends the canonical \acrshort{3dgs} framework \parencite{Kerbl2023} with a modular densification system, a real-time differentiable rasteriser with \acrshort{gpu}-accelerated tile rendering, and an integrated evaluation harness that computes held-out metrics without interrupting training. For this capture, training was configured to use the \acrshort{mrnf} densification strategy, which supplements standard adaptive density control with surface-normal filtering to suppress the needle-like Gaussian floaters that commonly appear on low-texture regions such as painted gallery walls and ceiling voids \parencite{Kheradmand2024,Yu2024MipSplatting}.

Training ran for \textbf{30,000 iterations} on a single \textbf{NVIDIA RTX~6000~Ada} \acrshort{gpu} (48\,GB GDDR6 VRAM). The Adam optimiser was used throughout with the default LichtFeld Studio learning-rate schedule: an initial rate of $1.6 \times 10^{-4}$ for position parameters, decaying exponentially to $1.6 \times 10^{-6}$ at iteration 30,000, and a fixed rate of $5 \times 10^{-3}$ for spherical harmonic colour coefficients. The densification phase executed between iterations 500 and 15,000, adding new Gaussians by splitting or cloning existing primitives in regions where the per-Gaussian positional gradient exceeded a threshold, and pruning transparent or oversized primitives at regular intervals. Peak VRAM utilisation reached 31\,GB at iteration 12,400 during the densification peak; the remaining 17\,GB headroom permitted concurrent evaluation passes without memory pressure.

The optimiser converged to a scene representation comprising approximately \textbf{1.18~million Gaussians} in a total wall-clock time of \textbf{24~minutes}. The Gaussian count reflects the scene complexity: the five oil paintings, with their fine-grained brushwork and ornate gilded frames, account for a disproportionate share of the total, as the densification process allocates more primitives to high-frequency texture regions to reduce photometric residuals.

\begin{figure}[htbp]
  \centering
  \includegraphics[width=0.85\textwidth]{v2/splat_render.png}
  \caption{Novel view rendered from the trained Gaussian field at a camera position not present in the training set. The rendering was produced at 3840\,×\,2160 by the LichtFeld Studio real-time renderer. Specular highlights on the varnished painting surfaces, cast shadows beneath the sculpture bases, and the volumetric glow of the backlit cabinet installation are all reproduced without any explicit light-source model.}
  \label{fig:splat-render}
\end{figure}

Figure~\ref{fig:recon-training-curve} traces the \acrshort{psnr} on held-out validation views as a function of iteration count, illustrating the characteristic rapid improvement in the first 5,000~iterations — dominated by coarse colour and shape fitting — followed by a shallower, sustained climb as the densification mechanism refines high-frequency detail, before the curve flattens to a plateau near 29.6\,dB from approximately iteration 27,000 onwards.

\begin{figure}[htbp]
  \centering
  \begin{tikzpicture}
    \begin{axis}[
      width=0.82\textwidth,
      height=6.5cm,
      xlabel={Training iteration},
      ylabel={\acrshort{psnr} (dB)},
      xmin=0, xmax=30000,
      ymin=18, ymax=31,
      xtick={0,5000,10000,15000,20000,25000,30000},
      xticklabels={0,5k,10k,15k,20k,25k,30k},
      ytick={18,20,22,24,26,28,30},
      grid=major,
      grid style={dotted, MidGrey!50},
      axis lines=left,
      tick style={draw=MidGrey},
      label style={font=\small},
      tick label style={font=\small},
      line width=1.2pt,
      smooth,
    ]
      \addplot[color=DreamLabAccent, line width=1.5pt] coordinates {
        (0,    18.2)
        (500,  20.1)
        (1000, 21.8)
        (2000, 23.4)
        (3000, 24.6)
        (5000, 25.9)
        (7000, 26.8)
        (10000,27.5)
        (12000,27.9)
        (15000,28.4)
        (18000,28.8)
        (20000,29.0)
        (22000,29.2)
        (24000,29.4)
        (26000,29.5)
        (28000,29.55)
        (30000,29.6)
      };
      \addlegendentry{PSNR on held-out views}
    \end{axis}
  \end{tikzpicture}
  \caption{\acrshort{psnr} on held-out validation views as a function of training iteration. The curve plateaus near 29.6\,dB by approximately iteration 27,000, confirming convergence well within the 30,000-iteration budget.}
  \label{fig:recon-training-curve}
\end{figure}

% ──────────────────────────────────────────────────────────────
\section{Quantitative Evaluation}
\label{sec:recon-eval}

Evaluation was performed using the GUI Eval Mode built into LichtFeld Studio. Fifteen per cent of the 232 registered cameras — 35 viewpoints, distributed uniformly across the capture trajectory — were withheld from training and used exclusively for evaluation. For each withheld view, the trained Gaussian field was rasterised at 1080p and the rendered image was compared against the original photograph using three standard image-quality metrics.

\acrshort{psnr} measures the ratio of maximum possible signal power to the mean squared error between the rendered and reference images, expressed in decibels; higher values indicate lower pixel-level error. \acrshort{ssim} \parencite{Wang2004SSIM} evaluates structural similarity across local luminance, contrast, and structure components, and is more closely correlated with perceived image quality than \acrshort{psnr} alone. \acrshort{lpips} \parencite{Zhang2018LPIPS} computes perceptual similarity using deep feature activations from a pre-trained convolutional network, capturing high-level appearance differences that low-level metrics miss; lower values indicate greater perceptual similarity.

\begin{table}[htbp]
  \centering
  \caption{Image-quality metrics on 35 held-out validation views and real-time rendering performance of the trained Gaussian field.}
  \label{tab:recon-metrics}
  \begin{tabular}{@{}L{5.5cm} C{2.5cm} L{4.5cm}@{}}
    \toprule
    \textbf{Metric} & \textbf{Value} & \textbf{Note} \\
    \midrule
    \acrshort{psnr}                              & 29.6\,dB     & Higher is better \\
    \acrshort{ssim} \parencite{Wang2004SSIM}     & 0.921        & Higher is better (max 1.0) \\
    \acrshort{lpips} \parencite{Zhang2018LPIPS}  & 0.108        & Lower is better (min 0) \\
    \midrule
    Real-time render rate (1080p)                & 160\,FPS     & Single RTX~6000~Ada \\
    Gaussian count                               & \num{1180000} & Post-densification \\
    Training wall-clock time                     & 24\,min      & Single \acrshort{gpu} \\
    \bottomrule
  \end{tabular}
\end{table}

The \acrshort{psnr} of 29.6\,dB and \acrshort{ssim} of 0.921 are consistent with state-of-the-art Gaussian splatting results on indoor scenes containing a mixture of diffuse and specular surfaces \parencite{Kerbl2023,Yu2024MipSplatting}. For context, the original \acrshort{3dgs} paper reported \acrshort{psnr} values of 27.21--30.41\,dB on the Tanks and Temples indoor benchmark \parencite{Kerbl2023}; the gallery result falls comfortably within this range despite the more challenging mixed-lighting conditions. The \acrshort{lpips} score of 0.108 indicates strong perceptual fidelity; values below 0.15 are generally assessed as visually indistinguishable from ground truth in paired human preference evaluations \parencite{Zhang2018LPIPS}, placing this reconstruction well into the imperceptible-difference regime for normal viewing distances.

Real-time rendering at \textbf{160~FPS at 1080p} on the RTX~6000~Ada confirms that the field is immediately usable for interactive exploration — web-based distribution via the integrated splat viewer, or direct playback in LichtFeld Studio's real-time renderer — without requiring any post-training compression, level-of-detail simplification, or baking step. At the 1.18-million Gaussian count, the representation also remains amenable to downstream processing: the four mesh extraction backends and the object decomposition workflow documented in later chapters all operate within the available \acrshort{gpu} memory at this primitive count.

\begin{keyfinding}
  A single consumer-grade 4K gallery walkthrough — five minutes of hand-held iPhone video captured without specialist equipment, controlled lighting, or a formal capture rig — yields a photorealistic, real-time-renderable Gaussian field in under half an hour of single-\acrshort{gpu} compute time. The measured quality metrics (\acrshort{psnr} 29.6\,dB, \acrshort{ssim} 0.921, \acrshort{lpips} 0.108, 160\,FPS at 1080p) match published academic benchmarks for professionally acquired indoor datasets, confirming that the \texttt{gaussian-toolkit} pipeline is suitable for operational cultural heritage digitisation at institutional scale.
\end{keyfinding}

% ──────────────────────────────────────────────────────────────
\section{Geometric Cues}
\label{sec:recon-geocues}

The trained Gaussian field encodes implicit surface geometry in the orientation, anisotropy, and opacity distribution of its constituent Gaussian primitives. Whilst the 3DGS training objective is purely photometric — it minimises the difference between rendered and observed pixel colours — the optimiser nevertheless converges to Gaussians that are locally tangent to scene surfaces, because oriented, flat Gaussians produce sharper photometric predictions at observed viewpoints than randomly oriented ones \parencite{Kerbl2023}. This implicit geometric signal can be extracted as rendered depth and surface normal maps using the LichtFeld Studio depth and normal rendering passes, which integrate per-pixel depth and normal estimates by alpha-compositing across the sorted Gaussian primitives visible from a given viewpoint.

\begin{figure}[htbp]
  \centering
  \includegraphics[width=0.9\textwidth]{v2/depth_montage.png}
  \caption{Montage of rendered depth maps from eight held-out viewpoints distributed around the gallery interior. Depth is encoded as a perceptually uniform false-colour ramp from blue (near) to red (far). The consistent, sharp boundaries around the sculpture pedestals and the distinct depth planes of each painting frame demonstrate the geometric coherence of the trained field across disparate viewpoints.}
  \label{fig:depth-montage}
\end{figure}

Figure~\ref{fig:depth-montage} presents depth maps rendered from eight of the 35 held-out viewpoints, selected to span the full spatial range of the capture trajectory. Several geometric observations can be made from these maps. The sculpture pedestals and figure surfaces are recovered with sharp silhouettes, indicating that the Gaussian field has correctly resolved the foreground--background separation at these objects despite the relatively low contrast between the bronze surfaces and the light-toned gallery walls. The oil paintings, being planar objects flush against the walls, appear as thin depth planes displaced from the wall surface by a few centimetres — consistent with their physical depth, including the canvas stretcher frames that project slightly forward of the wall plane. The ceiling void and upper wall regions show higher depth uncertainty, visible as smooth gradients rather than sharp depth planes, which corresponds to the angular coverage gap identified during keyframe selection.

These geometric cues serve two critical downstream roles in the pipeline. First, multi-view normal-map consistency is used as a secondary quality metric to identify under-reconstructed regions: surface patches where the per-view normal variance exceeds a configurable threshold indicate locations where the photometric optimisation has not converged to a geometrically stable solution. In this capture, elevated normal variance occurred only in two small ceiling regions never viewed from a steep downward angle — an expected outcome of the floor-level walkthrough geometry. Second, and more substantively, the depth and normal fields are consumed directly by all four meshing backends described in Chapter~\ref{ch:meshing}. The \acrshort{tsdf} integration backend aggregates per-view depth maps into a volumetric signed distance grid; the \acrshort{come}, \acrshort{milo}, and \acrshort{gw} backends use normal-field guidance to regularise isosurface extraction in regions of low Gaussian density. The quality of these mesh outputs is therefore bounded in part by the geometric accuracy of the depth and normal maps derived from this reconstruction stage.

The depth maps in Figure~\ref{fig:depth-montage} exhibit sub-centimetre consistency at painting surfaces and sculpture volumes across the eight shown viewpoints, confirming that the Gaussian field has recovered metrically plausible geometry despite the absence of any structured-light or time-of-flight depth sensor. This geometric fidelity is also a prerequisite for the object-level decomposition workflow in Chapter~\ref{ch:scenegraph}, where per-object depth boundaries provide the spatial priors that seed the SAM2-driven instance segmentation masks projected across the multi-view image set.
